\begingroup
\small
\setlength{\tabcolsep}{3pt}
\begin{longtable}{@{}L{0.08\textwidth}L{0.16\textwidth}L{0.69\textwidth}@{}}
\caption{The sixteen portable Cell48 opcodes}\label{tab:opcodes}\\
\toprule
\textbf{Code} & \textbf{Name} & \textbf{State effect} \\
\midrule
\endfirsthead
\multicolumn{3}{l}{\small\itshape The sixteen portable Cell48 opcodes -- continued}\\
\toprule
\textbf{Code} & \textbf{Name} & \textbf{State effect} \\
\midrule
\endhead
\midrule
\multicolumn{3}{r}{\small continued on next page}\\
\endfoot
\bottomrule
\endlastfoot
0 & NOP & No state mutation beyond lineage and successor selection. \\
1 & SET & Set one State64 field selected by flags[3:0] to arg0. \\
2 & JIT1 & Derive a deterministic parity bit and add +/\allowbreak{}-arg0 to the selected field. \\
3 & KIN2 & Apply integer second-order update v'=v+args;\allowbreak{} q'=q+v'. \\
4 & PHI & Advance lower-case phi modulo 2\textasciicircum{}12 and expose wrap/\allowbreak{}half-circle branch. \\
5 & TIME & Advance zero-based modular tick modulo 2\textasciicircum{}14 and fold winding into lineage. \\
6 & SDF0 & Return literal zero relation for the current defined LUT cell. \\
7 & CONE & Evaluate a quantized log-polar cone/\allowbreak{}support relation. \\
8 & SPHERE & Evaluate a quantized radial/\allowbreak{}phi shell relation. \\
9 & KLEIN & Apply reflective Klein or source half-turn radial wrap. \\
10 & RADIX & Select one bit of the canonical 64-bit key. \\
11 & HINGE & Apply map1 deltas and optional orientation/\allowbreak{}sheet flips when branch=1. \\
12 & LSYS & Apply chirality-signed lower-case-phi turn and dyadic velocity scaling. \\
13 & PROJECT & Set a symbolic downstream output token without halting. \\
14 & EMIT & Set output token and emitted status;\allowbreak{} flag bit 0 may terminate the program. \\
15 & HALT & Terminate the current state execution. \\
\end{longtable}
\normalsize
\endgroup
